\newpage
\pagestyle{fancy}

\section{Riferimenti legislativi e normativi}

Vengono di seguito riportate le principali norme e leggi applicabili. Ogni altra prescrizione, regolamentazione e raccomandazione emanata da eventuali enti ed applicabili agli impianti elettrici ed alle loro parti componenti deve essere recepita e attuata.

\begin{longtable}{|p{4.2cm}|p{9.5cm}|}

\hline
D.lgs. 09/04/08 n. 81 & ``Attuazione dell'art. 1 della Legge 3 agosto 2007, n. 123, in materia di tutela della salute e della sicurezza nei luoghi di lavoro.'' \\
\hline
Decreto 22/01/2008 n. 37 & ``Regolamento concernente l'attuazione dell'art.11-quaterdecies comma 13, lettera a) della Legge n. 248 del 2 dicembre 2005, recante riordino delle disposizioni in materia di attività di installazione degli impianti all'interno degli edifici.'' \\
\hline
D.P.R. 22/10/01, n. 462 & ``Regolamento di semplificazione del procedimento per la denuncia di installazioni e dispositivi di protezione contro le scariche atmosferiche, di dispositivi di messa a terra di impianti elettrici e di impianti elettrici pericolosi.'' \\
\hline
Legge 18/10/1977, n. 791 & ``Attuazione della direttiva del consiglio della Comunità Europea (n.73/73/CEE) relativa alle garanzie di sicurezza che deve possedere il materiale elettrico destinato ad essere utilizzato entro alcuni limiti di tensione.'' \\
\hline
Legge 01/03/1968, n. 186 & ``Disposizioni concernenti la produzione di materiali, apparecchiature, macchinari, installazione impianti elettrici ed elettronici.'' \\
\hline
Legge 05/03/1990, n. 46 & ``Disposizioni concernenti gli impianti elettrici, idrici ed idrosanitari radiotelevisivi ed elettronici in genere, nonché gli impianti per il trasporto e l'utilizzazione del gas e per gli impianti di riscaldamento e climatizzazione per regolamentare la sicurezza degli impianti negli edifici.'' \\
\hline
Regolamento prodotti da costruzione (UE) 305/2011 -- CPR & ``Regolamento (UE) N. 305/2011 del Parlamento Europeo e del Consiglio del 9 marzo 2011 che fissa condizioni armonizzate per la commercializzazione dei prodotti da costruzione e che abroga la direttiva 89/106/CEE del Consiglio (Testo rilevante ai fini del SEE).'' \\
\hline
D.lgs. 16/06/17 n. 106 & ``Adeguamento della normativa nazionale alle disposizioni del regolamento (UE) n.305/2011, che fissa condizioni armonizzate per la commercializzazione dei prodotti da costruzione e che abroga la direttiva 89/106/CEE''. \\
\hline
D.P.R. 06/12/1991, n. 447 & ``Disposizioni concernenti l'attuazione della Legge 05/03/1990, n. 46 in materia di sicurezza degli impianti''. \\
\hline
D.P.R. 18/04/1994, n. 392 & ``Regolamento recante disciplina del procedimento di riconoscimento delle imprese ai fini dell'installazione, ampliamento e trasformazioni degli impianti nel rispetto delle norme di sicurezza''. \\
\hline
D.P.R. 01/08/11 n. 151 & ``Regolamento recante semplificazione della disciplina dei procedimenti relativi alla prevenzione degli incendi, a norma dell'articolo 49, comma 4-quater, del decreto-legge 31 maggio 2010, n. 78, convertito, con modificazioni, dalla legge 30 luglio 2010, n. 122''. \\
\hline
\end{longtable}


\hspace{1em}
Norme del Comitato Elettrotecnico Italiano e Norme Europee relative alla esecuzione degli impianti richiesti, in quanto applicabili; in particolare sono inderogabili le Norme CEI seguenti:

\begin{longtable}{|p{4cm}|p{10cm}|}
\hline
CEI 0-2 & ``Guida per la definizione della documentazione di progetto degli impianti elettrici, elettronici e di comunicazione elettronica (EEC)''. \\
\hline
CEI 64-8 & ``Impianti elettrici utilizzatori a tensione non superiore a 1000V in C.A. e 1500V in C.C.'' \\
\hline
CEI 64-50 & ``Guida per l'integrazione nell'edificio degli impianti elettrici utilizzatori, ausiliari e telefonici''. \\
\hline
CEI 64-12 & ``Guida per l'esecuzione dell'impianto di terra negli edifici per uso residenziale''. \\
\hline
CEI 11-17 & ``Impianti di produzione, trasporto e distribuzione di energia elettrica. Linee in cavo''. \\
\hline
CEI 23-51 & ``Prescrizioni per la realizzazione, le verifiche e le prove dei quadri di distribuzione per installazioni fisse per uso domestico e similare''. \\
\hline
CEI 23-48 & ``Involucri per apparecchi per installazioni elettriche fisse per usi domestici e similari. Parte 1: prescrizioni generali''. \\
\hline
CEI 23-49 & ``Involucri per apparecchi per installazioni elettriche fisse per usi domestici e similari. Parte 2: prescrizioni particolari per involucri destinati a contenere dispositivi di protezione ed apparecchi che nell'uso ordinario dissipano una potenza non trascurabile''. \\
\hline
CEI 17-5 & ``Apparecchiature a bassa tensione -- interruttori automatici''. \\
\hline
CEI 17-11 & ``Interruttori di manovra, sezionatori, interruttori di manovra sezionatori e unità combinate con fusibili''. \\
\hline
CEI 17-44 & ``Apparecchiature a bassa tensione -- regole generali''. \\
\hline
CEI 17-113 & ``Apparecchiature assiemate di protezione e di manovra per bassa tensione (quadri bt). Parte 1: Regole generali''. \\
\hline
CEI 17-114 & ``Apparecchiature assiemate di protezione e di manovra per bassa tensione (quadri bt). Parte 2: Quadri di potenza''. \\
\hline
CEI 0-10 & ``Guida alla manutenzione degli impianti elettrici''. \\
\hline
CEI 0-11 & ``Guida alla gestione in qualità delle misure per la verifica degli impianti elettrici ai fini della sicurezza''. \\
\hline
CEI 81-10/1 & ``Protezione contro i fulmini. Principi generali''. \\
\hline
CEI 81-10/2 & ``Protezione contro i fulmini. Valutazione del rischio''. \\
\hline
CEI 81-10/3 & ``Protezione contro i fulmini. Parte 3: danno materiale alle strutture e pericolo per le persone''. \\
\hline
CEI 81-10/4 & ``Protezione contro i fulmini. Impianti elettrici ed elettronici inerenti alle strutture''. \\
\hline
CEI UNEL 35026 & ``Cavi elettrici isolati con materiale elastomerico e termoplastico per tensioni nominali di 1000 V in corrente alternata e 1500 V in corrente continua. Portate di corrente in regime permanente per posa interrata''. \\
\hline
CEI UNEL 35024/1 & ``Cavi elettrici isolati con materiale elastomerico e termoplastico per tensioni nominali di 1000 V in corrente alternata e 1500 V in corrente continua. Portate di corrente in regime permanente per posa in aria''. \\
\hline
CEI UNEL 35023 & ``Cavi per energia isolati in gomma o con materiale termoplastico aventi grado di isolamento non superiore a 4. Cadute di tensione''. \\
\hline
CEI UNEL 35016 & ``Classi di reazione al fuoco dei cavi elettrici in relazione al regolamento UE prodotti da costruzione (305/2011).'' \\
\hline
CEI 3-50 & ``Segni grafici da utilizzare sulle apparecchiature. Parte 2: segni originali''. \\
\hline
CEI 64-100/1 & ``Edilizia residenziale. Guida per la predisposizione delle infrastrutture per gli impianti elettrici, elettronici e per le comunicazioni. Parte 1: Montanti degli edifici''. \\
\hline
CEI 64-100/2 & ``Edilizia residenziale. Guida per la predisposizione delle infrastrutture per gli impianti elettrici, elettronici e per le comunicazioni. Parte 2: Unità immobiliari (appartamenti)''. \\
\hline
CEI 64-14 & ``Guide alle verifiche degli impianti elettrici utilizzatori''. \\
\hline
CEI 64-17 & ``Guida all'esecuzione degli impianti elettrici nei cantieri''. \\
\hline
CEI 64-53 & ``Edilizia residenziale. Guida per l'integrazione nell'edificio degli impianti elettrici utilizzatori e per la predisposizione di impianti ausiliari, telefonici e di trasmissione dati. Criteri particolari per edifici ad uso prevalentemente residenziale''. \\
\hline
CEI EN 60529 (CEI 70-1) & ``Gradi di protezione degli involucri (Codice IP)''. \\
\hline
CEI UNEL 35011 & ``Cavi per energia e segnalamento. Sigle di designazione.'' -- Cavi nazionali. \\
\hline
CEI UNEL 00722 & ``Identificazione delle anime dei cavi''. \\
\hline
CEI UNEL 00712 & ``Cavi, cordoni e fili per telecomunicazioni a bassa frequenza. Colori distintivi dell'isolante''. \\
\hline
CEI 23-50 & ``Spine e prese per usi domestici e similari. Parte 1: prescrizioni generali''. \\
\hline
CEI 11-25 & ``Correnti di cortocircuito nei sistemi trifase in corrente alternata. Parte 0: calcolo delle correnti''. \\
\hline
EN 50525-1/A1 CEI 20-107 & ``Cavi energia con tensione nominale non superiore a 450/750 V. Parte 1: prescrizioni generali''. \\
\hline
CEI 20-40/1-1 & ``Cavi elettrici. Guida all'uso dei cavi con tensione nominale non superiore a 450/750 V (U0/U) -- Parte 1: criteri generali.'' \\
\hline
CEI 20-40/2-1;V1 & ``Cavi elettrici. Guida all'uso dei cavi con tensione nominale non superiore a 450/750 V (U0/U) -- Parte 2: Criteri specifici relativi ai tipi di cavo specificati nella Norma EN 50525.'' \\
\hline
CEI 20-115 (EN 50575) & ``Cavi per energia, controllo e comunicazioni. Cavi per applicazioni generali nei lavori di costruzione soggetti a prescrizioni di resistenza all'incendio''. \\
\hline
CEI 20-35/1-2 (EN 60332-1-2) & ``Cavi non propaganti la fiamma''. \\
\hline
CEI 20-37/2 & ``Acidità e conduttività dei gas emessi durante la combustione dei cavi''. \\
\hline
CEI 20-37/3-1 (EN 61034-2) & ``Trasmittanza del fumo da parte dei cavi''. \\
\hline
CEI 20-21 & ``Calcolo delle portate dei cavi elettrici in regime permanente''. \\
\hline
CEI 0-16 & ``Regola tecnica di riferimento per la connessione di Utenti attivi e passivi alle reti AT ed MT delle imprese distributrici di energia elettrica''. \\
\hline
CEI 0-21 & ``Regola tecnica di riferimento per la connessione di Utenti attivi e passivi alle reti BT delle imprese distributrici di energia elettrica''. \\
\hline
\end{longtable}

Inoltre, tutte le apparecchiature facente parte dell’impianto oggetto del seguente intervento devono rispettare tutte le leggi, le norme vigenti in materia, anche se non citate espressamente e non richiamate, e le prescrizioni delle autorità locali, VV.F, ente distributore di energia elettrica, ecc.